\chapter{Systems of linear equations}
%\chapterepigraph{An optional short quotation or guiding question.}{Attribution}

\introsection
Systems of linear equations represent the very first sets of equations that we learn to solve. As we go deeper into their understanding, we begin to unleash the power of linear algebra, only paired by its beauty. 

\section{Solving systems of linear equations}
\label{sec:solvingsystems}
\sidenote{Gauss, Jordan}
A linear equation is an expression of the form $a_{1}x_{1}+...+a_{n}x_{n} = b$, where the coefficients $a_{i}$ and the independent term $b$ are \textbf{scalars}, that is, elements of the \textbf{field} $\R$ or $\C$, that will be specified beforehand, and the $x_{i} $ are indeterminates or variables. 

The equation $0=0$ is called \emph{trivial} and the equation $0=b$ with $b\neq 0$ is called \emph{inconsistent}. Any other equation is called \emph{consistent}.

\begin{definition}
  A  \textbf{system of linear equations} or linear system is a finite list of linear equations:
  \begin{equation*}
    (S) \linsys{a}{x}{b}{m}{n}
  \end{equation*}
\end{definition}

\begin{definition}
  A \textbf{vector} is an ordered tuple of scalars (numbers) of the field. Any of the scalars of a vector $v$ is called a \emph{coordinate} and if the vector has $n$ coordinates that are, for example, real numbers, we say that $v\in \R^{n}$. We will write vectors horizontally like
  \begin{equation*}
      v = \begin{pmatrix}
        v_{1} \\
        \vdots \\
        v_{n}
      \end{pmatrix}.
  \end{equation*}
  Nevertheless, we will usually write $v = (v_{1},\ldots,v_{n})^T$ to save space, meaning the same vector written horizontally like the one before. Later on, we will see that this $v^T$ operation has a name and meaning. \end{definition}

A \textbf{solution} to a linear system like (S) is a vector $(c_{1},\ldots,c_{n})$ such that, when we substitute the indeterminates $x_1,...,x_n$ by $c_1,...,c_n$ respectively, all the equations are satified. In other words,
\begin{equation*}
  \linsys{a}{c}{b}{}{}
\end{equation*}

If we look from a solutions perspective, we can classify equations in three types:
\begin{enumerate}
  \item The equation $0=0$ which is always true, and we call \textbf{trivial}.
  \item The equation $0=b$ with $b\neq 0$, which is always false and we call \textbf{inconsistent}.
  \item Any other equation has one solution, and we call them \textbf{consistent}.
\end{enumerate}

In the same manner, we can classify linear systems by their solutions:
\begin{enumerate}
  \item If every vector is a solution of $(S)$, we call it \textbf{trivial}. Note that this is only the case when all equations are trivial.
  \item If no vector is a solution of $(S)$, we call it \textbf{inconsistent}. This can be the case without any inconsistent equations.
  \item In any other case (the system has solutions, but not every vector is a solution), we call $(S)$ \textbf{consistent}:
    \subitem If there is only one solution, we call $(S)$ \textbf{determined}.
    \subitem If there is more than one solution, we call $(S)$ \textbf{underdetermined}.
\end{enumerate}

From the perspective of the term $b$ of the system, sometimes called the vector of \textbf{independent terms}, we can speak about two types of systems:
\begin{enumerate}
  \item \textbf{Homogeneus} systems, where $b=0$, so all equations end in $...=0$.
  \item \textbf{Nonhomogeneus} systems, where $b\neq 0$, so at least one equation does not end in $...=0$.
\end{enumerate}

Note that, to every nonhomogeneus system like $(S)$ we can associate a homogeneus system by replacing $b$ by zeros:
\begin{equation*}
  (S) \linsys{a}{x}{b}{}{} \longrightarrow (H) \linsys{a}{x}{0}{}{}
\end{equation*}

We observe that, if we take two solutions of a system $(S)$, call them $c=\hvc{c}{}^T$ and $d=\hvc{d}{}^T$, then their difference $c-d$ is a solution of the homogeneus system $(H)$ associated to $(S)$: indeed, to see that $c-d=(c_{1}-d_{1},...,c_{n}-d_{n})$ is a solution of $(H)$ we have to check that, when we substitute every $x_{i}$ by $c_{i}-d_{i}$ in $(H)$,
\begin{gather*}
  a_{j1}(c_{1}-d_{1})+...+a_{j}(c_{n}-d_{n}) = a_{1j}c_{1}-a_{1j}d_{1}+...+a_{jn}c_{n}-a_{jn}d_{n} =\\
  =a_{j1}c_{1}+...+a_{jn}c_{n}-(a_{j1}d_{1}-a_{jn}d_{n}) = b_{j}-b_{j}=0.
\end{gather*}
This proves that, when we substitute $x_{i}$ by $c_{i}-d_{i}$ in every equation $a_{j1}x_{1}+...+a_{jn}c_{n}$, we get $0$, so $c-d$ is a solution of $(H)$.

In the same way, every solution for $(S)$ can be obtained taking just one solution of $(S)$ and adding solutions for the associated homogeneus system $(H)$: if $c$ is a solution for $(S)$, and we want to see that every other solution $d$ of $(S)$ can be obtained as $c + \text{ solution for } (H)$, take $d = c+(c-d)$. By the previous remark, $c-d$ is a solution for $(H)$.

\begin{definition}
    Two systems are called equivalent if they have the same solutions. Formally, if we call the systems $(S_1)$ and $(S_2)$, this is to say
    \begin{equation*}
      \hvc{c}{} \text{ is a solution for } (S_1) \iff \hvc{c}{} \text{ is a solution for } (S_2)
    \end{equation*}
\end{definition}

In order to solve linear systems, we will operate with their equations to simplify them trying not to lose solutions in the process. What are the operations that allow us to transform systems in this way? They are called elementary operations:



\begin{textbox}{Key idea}
This is an optional framed box. Use \texttt{textblock} for an unframed
block or \texttt{shadedbox} for a pale gray background.
\end{textbox}
\begin{shadedbox}{Key idea}
  Shaded box
\end{shadedbox}

\begin{exercises}
\begin{exercise}\label{ex:first}
Prove that \(T(-x)=-T(x)\) for every linear map \(T\).
\end{exercise}
\end{exercises}
